% ============================================================
%  Chapter 2 — Fields, Finite Fields and Modular Arithmetic
% ============================================================
\section{Fields, Finite Fields and Modular Arithmetic}

A \emph{field} is the algebraic structure that provides the scalars for a vector
space.  Understanding the field axioms carefully—and seeing them satisfied in
finite and modular settings—builds the foundation for all of linear algebra.

% ------------------------------------------------------------
\begin{defbox}[Definition --- Field Axioms]
A \textbf{field} is a set $\F$ equipped with two binary operations, addition
$+$ and multiplication $\cdot$, satisfying the following axioms for all
$x,y,z\in\F$:

\medskip
\textbf{Additive axioms}
\begin{enumerate}[label=A\arabic*.]
    \item \textit{Commutativity:} $x+y=y+x$.
    \item \textit{Associativity:} $x+(y+z)=(x+y)+z$.
    \item \textit{Identity:} there exists a unique element $0\in\F$ s.t.\ $x+0=x$.
    \item \textit{Inverse:} for each $x$ there exists a unique element $-x\in\F$
          s.t.\ $x+(-x)=0$.
\end{enumerate}

\medskip
\textbf{Multiplicative axioms}
\begin{enumerate}[label=M\arabic*.]
    \item \textit{Commutativity:} $xy=yx$.
    \item \textit{Associativity:} $x(yz)=(xy)z$.
    \item \textit{Identity:} there exists a unique element $1\neq 0$ in $\F$
          s.t.\ $x\cdot 1=x$.
    \item \textit{Inverse:} for each $x\neq 0$ there exists a unique $x^{-1}\in\F$
          s.t.\ $xx^{-1}=1$.
\end{enumerate}

\medskip
\textbf{Distributivity}
\begin{enumerate}[label=D.]
    \item $x(y+z)=xy+xz$.
\end{enumerate}
\end{defbox}

% ------------------------------------------------------------
\begin{defbox}[Definition --- Finite Fields (Galois Fields)]
A \textbf{finite field} (also called a \emph{Galois field}) is a field with
finitely many elements.  The most common examples are the sets of integers
modulo a prime $p$, denoted $\F_p = \Z_p = \{0,1,\ldots,p-1\}$ with addition
and multiplication performed modulo $p$.
\end{defbox}

% ------------------------------------------------------------
\begin{defbox}[Definition --- Modular Arithmetic]
For integers $a,b$ and a positive integer $n$, we write
\[
    a \equiv b \pmod{n}
\]
to mean $n \mid (a-b)$, i.e.\ $a - b = kn$ for some integer $k$.
\end{defbox}

% ------------------------------------------------------------
\begin{tipbox}[Remark --- Intuition for Modular Arithmetic]
Modular arithmetic is arithmetic where integers ``wrap around'' at the modulus.
For example, $17 \equiv 5 \pmod{12}$ because $17-5=12$, which is divisible by
$12$.  Clocks are a familiar everyday instance: $13\text{ o'clock} \equiv 1
\pmod{12}$.
\end{tipbox}

% ------------------------------------------------------------
\begin{exbox}[Example --- The Field $\F_3$]
Consider $\F_3 = \{0,1,2\}$ with addition and multiplication modulo $3$.

\medskip
\textbf{Addition table ($\bmod\ 3$):}
\[
\begin{array}{c|ccc}
+ & 0 & 1 & 2 \\\hline
0 & 0 & 1 & 2 \\
1 & 1 & 2 & 0 \\
2 & 2 & 0 & 1
\end{array}
\]

\medskip
\textbf{Multiplication table ($\bmod\ 3$):}
\[
\begin{array}{c|ccc}
\cdot & 0 & 1 & 2 \\\hline
0     & 0 & 0 & 0 \\
1     & 0 & 1 & 2 \\
2     & 0 & 2 & 1
\end{array}
\]

All field axioms can be verified by inspection of these tables (proof omitted).
\end{exbox}

% ------------------------------------------------------------
\begin{thmbox}[Theorem --- $\Z_p$ is a Field for Prime $p$]
Let $p$ be a prime.  Then $\Z_p = \{0,1,\ldots,p-1\}$ with addition and
multiplication modulo $p$ is a field with $p$ elements.

\medskip
The field axioms (commutativity, associativity, identity, inverse,
distributivity) can each be verified; in particular, every nonzero element
$a \in \Z_p$ has a multiplicative inverse (since $\gcd(a,p)=1$ when $p$ is
prime and $0<a<p$).
\end{thmbox}

% ------------------------------------------------------------
\begin{warnbox}[Notation Notice --- $\Z_p$ vs.\ $\F_p$]
Both $\Z_p$ and $\F_p$ are used in these notes to denote the finite field of
integers modulo a prime $p$.  Some texts reserve $\Z_p$ for the $p$-adic
integers (a very different object) and use only $\F_p$.
In these notes the two notations are interchangeable and always mean
$\{0,1,\ldots,p-1\}$ with arithmetic modulo $p$.
\end{warnbox}
